% Part: computability
% Chapter: recursive-functions
% Section: primes

\documentclass[../../../include/open-logic-section]{subfiles}

\begin{document}

\olfileid{cmp}{rec}{pri}
\olsection{اعداد اول}

سوربندی کراندار و کمینه‌سازی کراندار ابزارهای فراوانی برای نشان‌دادن
بازگشتی اولیه بودن توابع و روابط طبیعی در اختیارمان می‌گذارند. برای
نمونه، رابطهٔ «$y$ بر~$x$ بخش‌پذیر است» را در نظر بگیرید که با
$x \mid y$ نوشته می‌شود. رابطهٔ $x \mid y$ هنگامی برقرار است که
تقسیم~$y$ بر~$x$ بی‌باقیمانده ممکن باشد؛ یعنی $y$ مضربی صحیح از~$x$
باشد. (اگر این رابطه برقرار نباشد، یعنی باقیماندهٔ تقسیم~$y$ بر~$x$
بزرگ‌تر از~$0$ باشد، می‌نویسیم $x \nmid y$.) به بیان دیگر،
$x \mid y$ هرگاه و تنها هرگاه برای یک~$z$، $x \cdot z = y$. آشکار
است که هر چنین~$z$ای، اگر وجود داشته باشد، باید $\leq y$ باشد. پس
$x \mid y$ هرگاه و تنها هرگاه برای یک $z \le y$،
$x \cdot z = y$. بنابراین می‌توان رابطهٔ $x \mid y$ را با سوربندی
وجودی کراندار از برابری و ضرب چنین تعریف کرد:
\[
x \mid y \defiff \bexists{z \leq y}{(x \cdot z) = y}.
\]
پس نشان داده‌ایم $x \mid y$ بازگشتی اولیه است.

عدد طبیعی~$x$ 
\emph{اول} است اگر نه~$0$ و نه~$1$ باشد و تنها بر~$1$ و خودش
بخش‌پذیر باشد. به بیان دیگر، اعداد اول چنان‌اند که هرگاه
$y \mid x$، یا $y=1$ یا $y=x$. برای آزمودن اول‌بودن~$x$، تنها باید
برای همهٔ $y \le x$ بررسی کنیم که آیا $y \mid x$؛ زیرا اگر $y>x$،
خودبه‌خود $y \nmid x$. پس رابطهٔ~$\fn{Prime}(x)$ که هرگاه و تنها
هرگاه $x$ اول باشد برقرار است، با
\[
\fn{Prime}(x) \defiff x \geq 2 \land \bforall{y \leq x}{(y \mid x \lif y
  = 1 \lor y = x)}
\]
تعریف می‌شود و بنابراین بازگشتی اولیه است.

اعداد اول عبارت‌اند از $2$، $3$، $5$، $7$، $11$ و جز آن. تابع~$p(x)$
را در نظر بگیرید که عدد اول~$x$ام در این دنباله را بازمی‌گرداند؛ یعنی
$p(0)=2$، $p(1)=3$، $p(2)=5$ و جز آن. (برای سهولت اغلب $p(x)$ را
به صورت~$p_x$ می‌نویسیم: $p_0=2$، $p_1=3$ و جز آن.)

اگر تابعی چون~$\fn{nextPrime(x)}$ داشته باشیم که نخستین عدد اول
بزرگ‌تر از~$x$ را بازگرداند، می‌توان~$p$ را به‌آسانی با بازگشت اولیه
تعریف کرد:
\begin{align*}
  p(0) & = 2\\
  p(x+1) & = \fn{nextPrime}(p(x)).
\end{align*}
چون $\fn{nextPrime}(x)$ کوچک‌ترین~$y$ای است که $y>x$ و~$y$ اول است،
می‌توان آن را به‌آسانی با جست‌وجوی نامحدود محاسبه کرد. اما به لطف
نتیجه‌ای از اقلیدس می‌توان آن را با کمینه‌سازی کراندار نیز تعریف کرد:
همواره عددی اول میان~$x$ و~$\fact{x}+1$ وجود دارد.
\[
  \fn{nextPrime(x)} =
  \bmin{y \leq \fact{x}+1}{(y > x \land \fn{Prime}(y))}.
\]
این نشان می‌دهد $\fn{nextPrime}(x)$ و در نتیجه~$p(x)$ نه‌تنها
محاسبه‌پذیر، بلکه بازگشتی اولیه‌اند.

(اگر کنجکاوید، این هم اثباتی کوتاه از قضیهٔ اقلیدس. فرض کنید $p_n$
بزرگ‌ترین عدد اول $\le x$ باشد و حاصل‌ضرب
$p=p_0\cdot p_1\cdot\dots\cdot p_n$ همهٔ اعداد اول~$\le x$ را در
نظر بگیرید. یا $p+1$ اول است یا عددی اول میان~$x$ و~$p+1$ وجود دارد.
چرا؟ فرض کنید $p+1$ اول نباشد. آنگاه عدد اولی چون~$q$ وجود دارد که
$q \mid p+1$ و $q<p+1$. هیچ‌یک از اعداد اول~$\le x$، عدد~$p+1$ را
تقسیم نمی‌کند. (بنا بر تعریف~$p$، هر عدد اول $p_i \le x$، عدد~$p$ را
بی‌باقیمانده تقسیم می‌کند. پس تقسیم~$p+1$ بر هر $p_i \le x$ باقیماندهٔ
$1$ دارد و بنابراین $p_i \nmid p+1$.) ازاین‌رو $q$ عددی اول است که
$>x$ و $<p+1$. همچنین $p \le \fact{x}$؛ پس عددی اول $>x$ و
$\le \fact{x}+1$ وجود دارد.)

\begin{prob}
تقسیم صحیح~$d(x,y)$ را با استفاده از کمینه‌سازی کراندار تعریف کنید.
\end{prob}

\end{document}
